\subsection{讨论}
\label{sec:discussion-zh}

\paragraph{主矩阵解读。}
表~\ref{tab:main-results-zh} 为固定延迟模型下的相对比较；物理重采下 B2 M2$=0.604$，B0 M2$=0.560$，M1 p50 由 374\,ms 降至 208\,ms。
B1 无闭环重采时 M2 与 B0 相同，说明重采策略是有效帧率提升的关键。

\paragraph{ROI 裁剪。}
COCO 舌区 ROI（8\% padding）使评分聚焦解剖区，但未消除 Laplacian 上的 clear/blur 重叠（表~\ref{tab:roi-ablation-zh}）；$\tau$ 由 0.498 变为 0.465。

\paragraph{质量注入。}
合成高斯模糊对齐验证协议；M5 Spearman $\rho{\approx}0.36$--$0.47$ 反映真实纹理重叠。

\paragraph{真实 B3。}
MobileNet 验证准确率 98.3\%；在 $\tau_{B2}$ 下 B3 M2$=0.934$，高于 B2 的 $0.604$，但 M1 p50 高约 137\,ms；$\tau_{B3}{=}0.5$ 的 B3t 与混合 B4 见表~\ref{tab:learned-hybrid-zh}。
主 claim 仍为可解释、低延迟的 B2；高 M2 场景可选 B3/B4 并单独标定 $\tau_{B3}$。

\paragraph{泛化与伦理。}
算法与 FR3 解耦；离线公开语料，JSONL 无患者标识。

\paragraph{要点汇总。}
\begin{enumerate}
  \item ShezhenV3 6719 张 + ROI 标定 $\tau{=}0.465$。
  \item B2 M2$=0.604$，B0 M2$=0.560$，B2 M1 p50 208\,ms vs.\ B0 374\,ms。
  \item B3/B4 M2$\approx 0.93$，需权衡 MobileNet 延迟与重试率。
  \item M6 辅轨不混入表 II。
\end{enumerate}

\paragraph{开放问题。}
混合 Edge-IQA+B3；端侧检测 ROI；FR3 运动模糊 augment。

\paragraph{临床转化。}
试点 JSONL + 本地 $\tau$ + 云端置信度对比。

\paragraph{Reviewer FAQ。}
\textit{为何不用纯深度 IQA？}延迟、flags、标定透明。
\textit{ROI 若分离度未升？}解剖聚焦 + 与 TCM 检测一致；B3 亦在 ROI 上训练。
\textit{B3 M2 更低仍报告？}全量矩阵显示学习型头在 $\tau_{B2}$ 下 M2 更高；需分指标讨论分类准确率与路由延迟。

\paragraph{带宽经济。}
边侧 $<$50\,ms 打分匹配 2--5\,fps 采集，避免 4K 全帧上云。

\paragraph{LangGraph vs ROS。}
图结构便于 Python ML 子进程与 rclpy 解耦时的日志导出。

\paragraph{仿真到真实差距。}
合成中位数 0.819/0.191 高估可分性；ShezhenV3 ROI 0.516/0.414 需报告重叠感知指标。

\paragraph{局限回顾。}
见 \S\ref{sec:failure-cases-zh}：合成模糊、负分离度、B3/$\tau$ 不匹配。

\paragraph{积极结论。}
在文档化延迟模型下，闭环 + 可解释 Edge-IQA 将 B0 有效上传率（M2$=0.560$）提升至 B2 的 $0.604$，并显著降低 RTT 中位数；学习型 B4 进一步将 M2 提至 $0.936$。

\paragraph{工程清单。}
固定 checkpoint SHA256；改 padding 必重标定。

\paragraph{数据可用性。}
划分与脚本开源；图像路径本地配置。

\paragraph{代码可用性。}
\texttt{edge\_iqa/}、\texttt{langgraph\_router/}、\texttt{franka\_api\_server/}。

\paragraph{Final remark.}
本文贡献是\emph{路由架构}在真实 TCM 图上的证据链，而非 IQA 回归榜冠军。
